\chapter{Arithmetic Functions}

\begin{notedefinition}[3.1]
An arithmetic function is a function $f:\N\to\C$.
\end{notedefinition}
Typical examples are
\[
 \tau(n)=\sum_{d\mid n}1,
 \qquad
 \sigma(n)=\sum_{d\mid n}d,
 \qquad
 \omega(n)=\#\{p:p\mid n,\ p\text{ prime}\},
\]
and $\pi(n)$, the number of primes not exceeding $n$.

\begin{notedefinition}[3.2]
A nonzero arithmetic function $f$ is multiplicative if
\[
 f(mn)=f(m)f(n)\qquad\text{whenever }\gcd(m,n)=1.
\]
It is completely multiplicative if this equality holds for all $m,n\in\N$.
\end{notedefinition}

\begin{notetheorem}[3.1]
If
\[
 n=p_1^{\alpha_1}\cdots p_r^{\alpha_r}
\]
is the prime factorization of $n$ and $f$ is multiplicative, then
\[
 f(n)=\prod_{i=1}^r f(p_i^{\alpha_i}).
\]
\end{notetheorem}
\begin{proof}
The prime powers are pairwise coprime, so repeated use of multiplicativity gives the formula.
\end{proof}

\begin{noteremark}
Every nonzero multiplicative function satisfies $f(1)=1$.  For fixed $k\in\C$, the function $n\mapsto n^k$ is completely multiplicative.  The functions $u(n)=1$ and $r(n)=n$ are the cases $k=0$ and $k=1$.  The prime-counting function $\pi$ is not multiplicative.
\end{noteremark}

For fixed $k\in\C$, define
\[
 \sigma_k(n)=\sum_{d\mid n}d^k.
\]

\begin{noteproposition}[3.1A]
The function $\sigma_k$ is multiplicative.
\end{noteproposition}
\begin{proof}
If $\gcd(m,n)=1$, every divisor $d$ of $mn$ has a unique representation $d=d_1d_2$ with $d_1\mid m$ and $d_2\mid n$.  Therefore
\[
 \sigma_k(mn)
 =\sum_{d_1\mid m}\sum_{d_2\mid n}(d_1d_2)^k
 =\sigma_k(m)\sigma_k(n).
\]
\end{proof}

\begin{notetheorem}[3.2]
If $n=p_1^{\alpha_1}\cdots p_r^{\alpha_r}$, then
\[
 \tau(n)=\prod_{i=1}^r(\alpha_i+1),
 \qquad
 \sigma(n)=\prod_{i=1}^r\frac{p_i^{\alpha_i+1}-1}{p_i-1}.
\]
\end{notetheorem}
\begin{proof}
For a prime power,
\[
 \tau(p^\alpha)=\alpha+1,
 \qquad
 \sigma(p^\alpha)=1+p+\cdots+p^\alpha
 =\frac{p^{\alpha+1}-1}{p-1}.
\]
Now apply Theorem 3.1.
\end{proof}

\begin{notetheorem}[3.3]
Euler's totient function is multiplicative.
\end{notetheorem}
\begin{proof}
For coprime $m,n$, the Chinese remainder theorem gives a bijection
\[
 (\Z/mn\Z)^\times\longrightarrow
 (\Z/m\Z)^\times\times(\Z/n\Z)^\times,
 \qquad [x]_{mn}\longmapsto([x]_m,[x]_n).
\]
Taking cardinalities gives $\totient(mn)=\totient(m)\totient(n)$.
\end{proof}

\begin{notecorollary}[3.4]
If $n=p_1^{\alpha_1}\cdots p_r^{\alpha_r}$, then
\[
 \totient(n)=\prod_{i=1}^r(p_i^{\alpha_i}-p_i^{\alpha_i-1})
 =n\prod_{i=1}^r\left(1-\frac1{p_i}\right).
\]
\end{notecorollary}

\begin{notetheorem}[3.5 (Gauss's totient sum)]
For every $n\in\N$,
\[
 \sum_{d\mid n}\totient(d)=n.
\]
\end{notetheorem}
\begin{proof}
Partition $\{1,2,\ldots,n\}$ according to the value of $\gcd(x,n)$.  For each divisor $d\mid n$, let
\[
 S_d=\{1\le x\le n:\gcd(x,n)=d\}.
\]
Writing $x=dy$ gives a bijection from $S_d$ to the reduced residue classes modulo $n/d$, so
$\abs*{S_d}=\totient(n/d)$.  Hence
\[
 n=\sum_{d\mid n}\abs*{S_d}
 =\sum_{d\mid n}\totient(n/d)
 =\sum_{d\mid n}\totient(d).
\]
\end{proof}

\begin{notedefinition}[3.3]
A positive integer $n$ is perfect if $\sigma(n)=2n$.
\end{notedefinition}

\begin{notetheorem}[3.6 (Euclid--Euler theorem)]
An even positive integer $n$ is perfect if and only if
\[
 n=2^{m-1}(2^m-1)
\]
for some $m\in\N$ such that $2^m-1$ is prime.
\end{notetheorem}
\begin{proof}
If $2^m-1$ is prime, multiplicativity gives
\[
 \sigma\bigl(2^{m-1}(2^m-1)\bigr)
 =(2^m-1)\cdot 2^m=2n.
\]

Conversely, write an even perfect number as $n=2^k n_1$, with $n_1$ odd.  Then
\[
 (2^{k+1}-1)\sigma(n_1)=2^{k+1}n_1.
\]
Since $\gcd(2^{k+1}-1,2^{k+1})=1$, write
$n_1=(2^{k+1}-1)n_2$.  The preceding equation becomes
\[
 \sigma(n_1)=2^{k+1}n_2.
\]
If $n_2>1$, then $1,n_2,n_1$ are distinct divisors of $n_1$, so
\[
 \sigma(n_1)\ge1+n_2+n_1=1+2^{k+1}n_2,
\]
a contradiction.  Hence $n_2=1$ and $n_1=2^{k+1}-1$.  Finally,
$\sigma(n_1)=2^{k+1}=n_1+1$, so $n_1$ has no positive divisors other than $1$ and itself and is prime.  Put $m=k+1$.
\end{proof}

\begin{notetheorem}[3.7]
If $m\in\N$ and $2^m-1$ is prime, then $m$ is prime.
\end{notetheorem}
\begin{proof}
If $m=ab$ with $a,b>1$, then
\[
 2^m-1=(2^a)^b-1=(2^a-1)(1+2^a+\cdots+2^{a(b-1)}),
\]
so $2^m-1$ is composite.
\end{proof}

\section{Dirichlet convolution}

\begin{notedefinition}[3.4]
For arithmetic functions $f$ and $g$, their Dirichlet convolution is
\[
 (f*g)(n)=\sum_{d\mid n}f(d)g(n/d).
\]
Let $I$ denote the identity function
\[
 I(n)=\begin{cases}1,&n=1,\\0,&n>1.\end{cases}
\]
\end{notedefinition}

\begin{notetheorem}[3.8]
For arithmetic functions $f,g,h$:
\begin{enumerate}[label=(\alph*)]
\item $(f*g)*h=f*(g*h)$;
\item $f*g=g*f$;
\item $f*I=f$.
\end{enumerate}
\end{notetheorem}
\begin{proof}
Commutativity follows by replacing $d$ with $n/d$.  For associativity,
\[
 ((f*g)*h)(n)
 =\sum_{abc=n}f(a)g(b)h(c)
 =(f*(g*h))(n).
\]
For the identity, only the divisor $d=n$ contributes to
$\sum_{d\mid n}f(d)I(n/d)$.
\end{proof}

\begin{notedefinition}[3.5]
If $f*g=I$, then $g$ is the Dirichlet inverse of $f$, denoted $f^{-1}$.
\end{notedefinition}

\begin{notetheorem}[3.9]
An arithmetic function $f$ has a Dirichlet inverse if and only if $f(1)\ne0$.
\end{notetheorem}
\begin{proof}
If $f*g=I$, then $f(1)g(1)=1$, so $f(1)\ne0$.

Conversely, if $f(1)\ne0$, define recursively
\[
 g(1)=\frac1{f(1)},
 \qquad
 g(n)=-\frac1{f(1)}
 \sum_{\substack{d\mid n\\d>1}}f(d)g(n/d)
 \quad(n>1).
\]
Then $(f*g)(1)=1$, and for $n>1$ the defining equation makes $(f*g)(n)=0$.
\end{proof}

\begin{notetheorem}[3.10]
If $f$ and $g$ are multiplicative, then $f*g$ is multiplicative.
\end{notetheorem}
\begin{proof}
Let $\gcd(m,n)=1$.  Every divisor of $mn$ is uniquely $d_1d_2$ with $d_1\mid m$ and $d_2\mid n$.  Hence
\begin{align*}
 (f*g)(mn)
 &=\sum_{d_1\mid m}\sum_{d_2\mid n}
 f(d_1d_2)g\left(\frac m{d_1}\frac n{d_2}\right)\\
 &=\left(\sum_{d_1\mid m}f(d_1)g(m/d_1)\right)
   \left(\sum_{d_2\mid n}f(d_2)g(n/d_2)\right).
\end{align*}
\end{proof}

\begin{noteremark}
The Dirichlet inverse of a multiplicative function is multiplicative.  Indeed, define a multiplicative function $h$ on prime powers by the recursion in Theorem 3.9.  Then $f*h=I$ on prime powers and, by multiplicativity, on every integer.  Uniqueness of the inverse gives $h=f^{-1}$.
\end{noteremark}

\begin{notetheorem}[3.11]
If $f$ is multiplicative and
\[
 F(n)=\sum_{d\mid n}f(d),
\]
then $F$ is multiplicative.
\end{notetheorem}
\begin{proof}
Since $F=f*u$ and both $f$ and $u$ are multiplicative, Theorem 3.10 applies.
\end{proof}
\begin{noteremark}
Taking $f(n)=n^k$ proves that $\sigma_k$ is multiplicative.  Taking $f=\totient$ and checking prime powers gives another proof of
$\sum_{d\mid n}\totient(d)=n$.
\end{noteremark}

\begin{notedefinition}[3.6]
The M\"obius function is
\[
 \mobius(n)=
 \begin{cases}
 0,&p^2\mid n\text{ for some prime }p,\\
 (-1)^{\omega(n)},&n\text{ is squarefree}.
 \end{cases}
\]
In particular, $\mobius(1)=1$.
\end{notedefinition}

\begin{notetheorem}[3.12]
The Dirichlet inverse of $u$ is $\mobius$; equivalently,
\[
 \sum_{d\mid n}\mobius(d)=I(n).
\]
\end{notetheorem}
\begin{proof}
Both sides are multiplicative, so it suffices to check prime powers.  For $n=p^\alpha$ with $\alpha\ge1$,
\[
 \sum_{d\mid p^\alpha}\mobius(d)=\mobius(1)+\mobius(p)=1-1=0.
\]
\end{proof}

\begin{notetheorem}[3.13 (M\"obius inversion)]
If
\[
 F(n)=\sum_{d\mid n}f(d),
\]
then
\[
 f(n)=\sum_{d\mid n}\mobius(d)F(n/d)
     =\sum_{d\mid n}\mobius(n/d)F(d).
\]
\end{notetheorem}
\begin{proof}
The hypothesis is $F=f*u$.  Convolving with $\mobius=u^{-1}$ gives $f=F*\mobius$.
\end{proof}

\begin{notetheorem}[3.14 (converse form)]
If $f=\mobius*F$, then $F=f*u$; in particular, the two formulas in Theorem 3.13 are equivalent.
\end{notetheorem}

\section{The zeta function and Dirichlet series}

\begin{notedefinition}[3.7]
For $s\in\C$ with $\Re(s)>1$, the Riemann zeta function is
\[
 \zeta(s)=\sum_{n=1}^{\infty}\frac1{n^s}.
\]
\end{notedefinition}

\begin{notetheorem}[3.15 (Euler product)]
For $\Re(s)>1$,
\[
 \zeta(s)=\prod_{p\ \mathrm{prime}}\frac1{1-p^{-s}}.
\]
\end{notetheorem}
\begin{proof}
For real $s>1$, let the finite product range over $p\le N$.  Expanding its geometric series gives the sum of $n^{-s}$ over positive integers all of whose prime factors are at most $N$.  This includes every $n\le N$, so
\[
 0\le \zeta(s)-\prod_{p\le N}(1-p^{-s})^{-1}
 \le\sum_{n>N}n^{-s}\longrightarrow0.
\]
For complex $s$ with $\Re(s)>1$, the same argument applies absolutely after replacing $s$ by $\Re(s)$ in the estimates.
\end{proof}

\begin{notedefinition}[3.8]
The Dirichlet series of an arithmetic function $f$ is
\[
 D(f,s)=\sum_{n=1}^{\infty}\frac{f(n)}{n^s}
\]
wherever the series converges.
\end{notedefinition}

\begin{notetheorem}[3.16]
If $D(f,s)$ and $D(g,s)$ converge absolutely, then $D(f*g,s)$ converges absolutely and
\[
 D(f,s)D(g,s)=D(f*g,s).
\]
\end{notetheorem}
\begin{proof}
Absolute convergence gives
\[
 \sum_{a,b\ge1}\left|\frac{f(a)g(b)}{(ab)^s}\right|<\infty,
\]
so the double series may be rearranged.  Grouping terms with $ab=n$ yields
\[
 D(f,s)D(g,s)
 =\sum_{n=1}^{\infty}\frac1{n^s}\sum_{ab=n}f(a)g(b)
 =D(f*g,s).
\]
\end{proof}

\begin{notetheorem}[3.17]
For $\Re(s)>1$,
\[
 \frac1{\zeta(s)}=\sum_{n=1}^{\infty}\frac{\mobius(n)}{n^s}.
\]
\end{notetheorem}
\begin{proof}
Since $u*\mobius=I$, Theorem 3.16 gives
\[
 D(u,s)D(\mobius,s)=D(I,s)=1.
\]
\end{proof}

\section{Two natural-density applications}

\begin{noteproposition}[3.18]
The set of squarefree positive integers has natural density
\[
 \frac1{\zeta(2)}=\frac6{\pi^2}.
\]
\end{noteproposition}
\begin{proof}
The indicator of squarefreeness is
\[
 \sum_{d^2\mid n}\mobius(d).
\]
If $Q(N)$ counts squarefree integers not exceeding $N$, then
\begin{align*}
 Q(N)
 &=\sum_{d\le\sqrt N}\mobius(d)\floor{N/d^2}\\
 &=N\sum_{d\le\sqrt N}\frac{\mobius(d)}{d^2}+O(\sqrt N).
\end{align*}
After division by $N$, the error tends to zero and the finite sum tends to
$\sum_{d\ge1}\mobius(d)/d^2=1/\zeta(2)$.
\end{proof}

\begin{noteproposition}[3.19]
The proportion of ordered pairs $(a,b)$ with $1\le a,b\le N$ and $\gcd(a,b)=1$ tends to
\[
 \frac1{\zeta(2)}=\frac6{\pi^2}.
\]
\end{noteproposition}
\begin{proof}
Using $I(\gcd(a,b))=\sum_{d\mid a,\ d\mid b}\mobius(d)$, the number $A(N)$ of coprime pairs is
\[
 A(N)=\sum_{d\le N}\mobius(d)\floor{N/d}^2.
\]
Since $\floor{N/d}^2=N^2/d^2+O(N/d)$,
\[
 \frac{A(N)}{N^2}
 =\sum_{d\le N}\frac{\mobius(d)}{d^2}
 +O\left(\frac1N\sum_{d\le N}\frac1d\right)
 \longrightarrow\frac1{\zeta(2)}.
\]
\end{proof}
